\section{Finite principal Rees refinements of actual open covers}

\begin{lemma}[A finite principal refinement generating the unit ideal]\label{lem:sc0-finite-principal-open-refinement}
\lean{EnvelopeProperty.Char0.exists_finite_principal_refinement}
\leanok
Every actual open cover of $\operatorname{Spec}R$ admits a finite principal-open
refinement $D(g_1),\ldots,D(g_n)$ such that $(g_1,\ldots,g_n)=R$.
\end{lemma}
\begin{proof}
At each point choose a covering open containing it and a basic open containing
the point and contained in that covering open. The spectrum is quasi-compact,
so finitely many of the chosen basic opens cover it. If their functions generated
a proper ideal, that ideal would be contained in a prime ideal, giving a point
outside all the chosen basic opens. Therefore the functions generate the unit
ideal. This also handles an empty spectrum, where the empty family generates the
unit ideal of the zero ring.
\end{proof}

\begin{corollary}[Actual finite Rees refinement subordinate to an open cover]\label{cor:sc0-finite-principal-rees-refinement}
\lean{EnvelopeProperty.Char0.exists_finite_principal_rees_refinement}
\leanok
\uses{lem:sc0-finite-principal-open-refinement,lem:sc0-principal-rees-finite-cover}
For a subalgebra $C\subseteq K$ of a field, every actual open cover of
$\operatorname{Spec}C$ admits a finite unit-generating Rees node whose actual
chart images are each contained in one covering open.
\end{corollary}
\begin{proof}
Apply the finite principal-refinement lemma and identify every actual nonzero
Rees chart image with the corresponding original basic open. The family has a
nonzero member by the actual principal Rees cover theorem, so it defines an
allowed Rees node.
\end{proof}
